\section{Question and hypothesis}
\label{sec:introduction}

Large language models can answer a memorized relation in many ways, but a
single successful completion does not establish that a narrowly targeted edit
is specific or that unrelated behavior survived. We therefore asked a bounded
question: can a small BF16 LoRA update teach the synthetic relation \Fact{} to
one pinned Qwen3.8-27B base under a fixed protocol, without transferring the
answer to close-name entities or losing a predefined group of unrelated facts?

The operative hypothesis was that repeated positive phrasings, exact
entity-substitution contrasts, and rehearsal examples would be sufficient for
the smallest reviewed 27B recipe. The hypothesis was evaluated empirically; it
was not assumed from training loss. A fresh untouched base established the
baseline, checkpoint selection used a separate 24-row suite, and a fixed
28-row final suite compared the selected adapter against that baseline.
\claimsource{evaluation}\claimsource{metadata}

The observed result met the registered policy. Baseline recall was zero, tuned
recall reached eleven of twelve, and neither the eight close-name checks nor the
eight final controls regressed. The repository's scorer consequently emitted
the cautious interpretation \Code{candidate-knowledge-acquisition}.
\claimsource{manifest} This is a description of one admitted result, not a claim
that the update mechanism was isolated, universally reliable, or uniquely
responsible for the behavioral change.

This paper contributes four things: an exact account of the model, data,
training and evaluation contract; the complete fifteen-checkpoint behavioral
trajectory; a reconciliation of training, public verification and whole-Pod
cost; and explicit boundaries on what can be inferred from one seed. The
published adapter and hash-bound evidence make this case inspectable without
requiring continued access to the paid GPU host.
